\begin{figure}[!htbp]
%	\centering
%	\def\ax{0.5}\def\ay{0.30}
%	\providecommand{\xyz}[3]{{#2 + (#1)*\ax},{#3 + (#1)*\ay}}
%	\providecommand{\rightangle}[3]{%
%		\coordinate (ra1) at ($(#2)!2.4mm!(#1)$);
%		\coordinate (ra3) at ($(#2)!2.4mm!(#3)$);
%		\draw[line width=0.5pt] (ra1) -- ($(ra1)+(ra3)-(#2)$) -- (ra3);}
%	\begin{minipage}[b]{0.46\textwidth}\centering
%		\begin{tikzpicture}[>=Stealth, line width=0.8pt, scale=1.0]
%			% ===== PANNELLO (a) =====
%			\draw[gray!60] (\xyz{-1.7}{-1.2}{0}) -- (\xyz{-1.7}{3.50}{0})-- (\xyz{3.0}{3.50}{0}) -- (\xyz{3.0}{-1.2}{0}) -- cycle;
%			\node[gray!75] at (\xyz{-1.2}{-1.0}{0}) {$\pi$};
%			\coordinate (O)  at (\xyz{0}{0}{0});
%			\coordinate (Pk) at (\xyz{0.5}{2.0}{0});
%			\coordinate (Ff) at (\xyz{1.7}{2.3}{0});
%			\coordinate (Of) at (\xyz{1.2}{0.3}{0});
%			\coordinate (MuB) at (\xyz{0}{0}{1.25});
%			\coordinate (MuC) at (\xyz{0}{0}{1.47});
%			\draw[gray!55, thin] (O) -- (Of) -- (Ff) -- (Pk) -- cycle;
%			
%			%				\draw[->, thin] (O) -- (Pk);
%			%				\node[font=\small, below] at ($(O)!0.5!(Pk)$) {$\overrightarrow{OP_k}$};
%			
%			\draw[->,thick] (O) -- (Pk) node[font=\small, midway, sloped, below] {$\overrightarrow{OP_k}$};
%			
%			\draw[->, thick] (Pk) -- (Ff) node[right=-2pt] {$\fb_k$};
%			\draw[->>, thick] (O) -- (MuB) node[left=1pt, yshift=-10pt] {$\bm{\upmu}_k(O)$};
%			%\draw[->, thin] ([shift={(0.18,0)}]MuC)arc[x radius=0.18, y radius=0.07, start angle=0, end angle=305];
%			\begin{scope}[gray!55]
%				\rightangle{MuB}{O}{Of}   % angolo retto fra mu e la retta // f_k
%				\rightangle{MuB}{O}{Pk}
%			\end{scope}
%			\fill (O)  circle (1.5pt) node[left=2pt] {$O$};
%			\fill (Pk) circle (1.5pt) node[below right=-1pt] {$P_k$};
%		\end{tikzpicture}\\[0.8ex]
%		(a)
%	\end{minipage}\hfill
%	\begin{minipage}[b]{0.50\textwidth}\centering
%		\begin{tikzpicture}[>=Stealth, line width=0.8pt, scale=1.0]
%			% proiezione: x in avanti (basso-sinistra), y destra, z su
%			\def\ax{-0.65}\def\ay{-0.48}
%			
%			\def\az{0.7}                            % <1 = comprime in altezza
%			\renewcommand{\xyz}[3]{{#2 + (#1)*\ax},{(#3)*\az + (#1)*\ay}}
%			% componenti di P_k
%			\def\xk{1.4}\def\yk{2.6}\def\zk{2.0}
%			\coordinate (O)   at (\xyz{0}{0}{0});
%			\coordinate (Px)  at (\xyz{\xk}{0}{0});
%			\coordinate (Py)  at (\xyz{0}{\yk}{0});
%			\coordinate (Pxy) at (\xyz{\xk}{\yk}{0});
%			\coordinate (Pk)  at (\xyz{\xk}{\yk}{\zk});
%			
%			% --- scatola di decomposizione (base + verticale) ---
%			\draw[gray!55, thin] (O) -- (Px) -- (Pxy) -- (Py) -- cycle;
%			\draw[gray!55, thin] (Pxy) -- (Pk);
%			
%			% --- assi x,y,z ---
%			\draw[->, thin] (O) -- (\xyz{2.6}{0}{0}) node[below left=-2pt] {$x$};
%			\draw[->, thin] (O) -- (\xyz{0}{3.4}{0}) node[right]          {$y$};
%			\draw[->, thin] (O) -- (\xyz{0}{0}{2.9}) node[above]          {$z$};
%			
%			% --- vettore posizione OP_k ---
%			%				\draw[->, thin] (O) -- (Pk);
%			%				\node[font=\small, above left=-8pt, yshift=6pt] at ($(O)!0.42!(Pk)$) {$\overrightarrow{OP_k}$};
%			
%			% --- terna delle componenti di forza in P_k ---
%			\coordinate (Xt) at (\xyz{\xk+0.85}{\yk}{\zk});
%			\coordinate (Yt) at (\xyz{\xk}{\yk+0.85}{\zk});
%			\coordinate (Zt) at (\xyz{\xk}{\yk}{\zk+0.85});
%			\draw[->, thick] (Pk) -- (Xt) node[above=4pt] {$X_k$};
%			\draw[->, thick] (Pk) -- (Yt) node[right=-1pt] {$Y_k$};
%			\draw[->, thick] (Pk) -- (Zt) node[above=-1pt] {$Z_k$};
%			
%			% --- etichette componenti (spigoli O->Py->Pxy->Pk) ---
%			\node[font=\small, below=20pt, xshift=-20pt]      at ($(O)!0.5!(Py)$)   {$y_k$};
%			\node[font=\small,  right=1pt, yshift=-6pt] at ($(Py)!0.5!(Pxy)$) {$x_k$};
%			\node[font=\small, right, yshift=6pt]      at ($(Pxy)!0.5!(Pk)$) {$z_k$};
%			
%			% --- punti ---
%			\fill (O)  circle (1.5pt) node[above left=-1pt]  {$O$};
%			\fill (Pk) circle (1.5pt) node[above right=-1pt] {$P_k$};
%			
%			% ===== MOMENTI in O (interni, assiali ->> + ricciolo) =====
%			\coordinate (Mx) at (\xyz{1.2}{0}{0});
%			\coordinate (My) at (\xyz{0}{1.2}{0});
%			\coordinate (Mz) at (\xyz{0}{0}{1.4});
%			\draw[->>, thick] (O) -- (Mx);
%			%\draw[->, thin, rotate around={126:(Mx)}]([shift={(0.16,0)}]Mx) arc[x radius=0.16, y radius=0.06, start angle=0, end angle=305];
%			\node[font=\small, above,  xshift=-4pt, yshift=6pt] at (Mx) {$\bm{\upmu}_{xk}(O)$};
%			\draw[->>, thick] (O) -- (My);
%			%\draw[->, thin, rotate around={-90:(My)}] ([shift={(0.16,0)}]My) arc[x radius=0.16, y radius=0.06, start angle=0, end angle=305];
%			\node[font=\small, below=-1pt, xshift=-6pt] at (My) {$\bm{\upmu}_{yk}(O)$};
%			\draw[->>, thick] (O) -- (Mz);
%			%\draw[->, thin, rotate around={0:(Mz)}]	([shift={(0.16,0)}]Mz) arc[x radius=0.16, y radius=0.06, start angle=0, end angle=305];
%			\node[font=\small, left=1pt] at (Mz) {$\bm{\upmu}_{zk}(O)$};
%			
%			% ===== FORZE lungo gli assi (esterne, -> ; traslabili con \sx,\sy,\sz) =====
%			\def\Lm{0.9}
%			\def\sx{1.2}
%			\coordinate (Xa) at (\xyz{\sx}{0}{0});
%			\coordinate (Xb) at (\xyz{\sx+\Lm}{0}{0});
%			\draw[->, thick] (Xa) -- (Xb);
%			\node[font=\small, below=3pt] at (\xyz{\sx+\Lm/2}{0}{0}) {$X_k$};
%			\def\sy{2.1}
%			\coordinate (Ya) at (\xyz{0}{\sy}{0});
%			\coordinate (Yb) at (\xyz{0}{\sy+\Lm}{0});
%			\draw[->, thick] (Ya) -- (Yb);
%			\node[font=\small, above] at (Yb) {$Y_k$};
%			\def\sz{1.5}
%			\coordinate (Za) at (\xyz{0}{0}{\sz});
%			\coordinate (Zb) at (\xyz{0}{0}{\sz+\Lm});
%			\draw[->, thick] (Za) -- (Zb);
%			\node[font=\small, left=2pt] at (Zb) {$Z_k$};
%			
%		\end{tikzpicture}
%		\\[0.8ex]
%		(b)
%	\end{minipage}
%	\caption{(a)~Momento $\bm{\upmu}_k(O)$ rispetto al polo
%		$O$ della forza $(P_k,\fb_k)$, definito come
%		$\bm{\upmu}^O_k = \protect\overrightarrow{OP_k}\times\fb_k$.
%		(b)~Componenti cartesiane $x_k,y_k,z_k$ del vettore
%		posizione $\protect\overrightarrow{OP_k}$ e componenti
%		$X_k,Y_k,Z_k$ della forza $\fb_k$.}
%	\label{fig:momento_forza}
%\end{figure}
